%=====================================================================
%  Trabalho 08 - Reconhecimento de Flores Iris usando MLP
%  EL056 - Redes Neurais Artificiais (PPGEELT/UFU)
%  Compilar dentro da pasta latex/:  pdflatex relatorio_trabalho08 (2x)
%  Resultados: marcadores \textbf{XX,X\%} substituidos automaticamente
%  por latex/resultados_macros.tex (gerado por gerar_macros_latex.py
%  a partir de resultados.json).
%=====================================================================
\documentclass[12pt,a4paper]{article}

%--------------------------- Pacotes ---------------------------------
\usepackage[utf8]{inputenc}
\usepackage[T1]{fontenc}
\IfFileExists{lmodern.sty}{\usepackage{lmodern}}{}
\usepackage[brazil]{babel}
\usepackage[top=3cm,left=3cm,bottom=2cm,right=2cm]{geometry}
\usepackage{setspace}
\usepackage{ragged2e}
\usepackage{graphicx}
\usepackage{amsmath,amssymb}
\usepackage{booktabs}
\usepackage{multirow}
\usepackage{float}
\usepackage{caption}
\usepackage{subcaption}
\usepackage{xcolor}
\usepackage{listings}
\usepackage{algorithm}
\usepackage{algpseudocode}
\usepackage{indentfirst}
\usepackage[hidelinks]{hyperref}

\onehalfspacing
\setlength{\parindent}{1.25cm}
\graphicspath{{../figs/}{figs/}}
\captionsetup{font=small,labelfont=bf,justification=centering}
\floatname{algorithm}{Algoritmo}
\renewcommand{\listalgorithmname}{Lista de algoritmos}
\algrenewcommand\algorithmicrequire{\textbf{Entrada:}}
\algrenewcommand\algorithmicensure{\textbf{Saída:}}

\lstset{
  language=Python, basicstyle=\ttfamily\footnotesize, keywordstyle=\color{blue!70!black}\bfseries,
  commentstyle=\color{green!45!black}\itshape, stringstyle=\color{red!60!black},
  numbers=left, numberstyle=\tiny\color{gray}, frame=single, breaklines=true,
  showstringspaces=false, columns=fullflexible, tabsize=4, captionpos=b,
  literate={á}{{\'a}}1 {ã}{{\~a}}1 {é}{{\'e}}1 {ç}{{\c{c}}}1 {ó}{{\'o}}1 {í}{{\'i}}1
}
\renewcommand{\lstlistingname}{Código}

%--------------------- Bloco de capa ABNT (PPGEELT/UFU) ---------------
\input{abnt_capa_1}

%=====================================================================
%  MARCADORES DE RESULTADOS (preenchidos via resultados_macros.tex)
%=====================================================================
\newcommand{\XX}{\textbf{XX,X\%}}
\newcommand{\XN}{\textbf{XX}}
\newcommand{\AccTeste}{\XX}      \newcommand{\FUmTeste}{\XX}     \newcommand{\CETeste}{\XN}
\newcommand{\AccVal}{\XX}        \newcommand{\AccTreino}{\XX}
\newcommand{\CETreino}{\XN}      \newcommand{\CEVal}{\XN}
\newcommand{\MelhorEpoca}{\XN}   \newcommand{\EpocasExec}{\XN}   \newcommand{\NParam}{\XN}
\newcommand{\GCSig}{\XN}         \newcommand{\GCTanh}{\XN}       \newcommand{\GCRelu}{\XN}
\newcommand{\PSet}{\XN}\newcommand{\RSet}{\XN}\newcommand{\FSet}{\XN}
\newcommand{\PVer}{\XN}\newcommand{\RVer}{\XN}\newcommand{\FVer}{\XN}
\newcommand{\PVir}{\XN}\newcommand{\RVir}{\XN}\newcommand{\FVir}{\XN}
\newcommand{\PMacro}{\XN}\newcommand{\RMacro}{\XN}\newcommand{\FMacro}{\XN}
\newcommand{\CMTAA}{\XN}\newcommand{\CMTAB}{\XN}\newcommand{\CMTAC}{\XN}
\newcommand{\CMTBA}{\XN}\newcommand{\CMTBB}{\XN}\newcommand{\CMTBC}{\XN}
\newcommand{\CMTCA}{\XN}\newcommand{\CMTCB}{\XN}\newcommand{\CMTCC}{\XN}
\newcommand{\CMKAA}{\XN}\newcommand{\CMKAB}{\XN}\newcommand{\CMKAC}{\XN}
\newcommand{\CMKBA}{\XN}\newcommand{\CMKBB}{\XN}\newcommand{\CMKBC}{\XN}
\newcommand{\CMKCA}{\XN}\newcommand{\CMKCB}{\XN}\newcommand{\CMKCC}{\XN}
\newcommand{\AccKfold}{\XX}      \newcommand{\DPKfold}{\XX}
\newcommand{\FUmKfold}{\XX}      \newcommand{\DPFUmKfold}{\XX}   \newcommand{\FoldsKfold}{\XX}
\newcommand{\ErrosVerVir}{\XN}   \newcommand{\ErrosTotalKfold}{\XN}
\newcommand{\AccTesteLR}{\XX}    \newcommand{\AccKfoldLR}{\XX}   \newcommand{\DPKfoldLR}{\XX}
\newcommand{\AccTesteKNN}{\XX}   \newcommand{\AccKfoldKNN}{\XX}  \newcommand{\DPKfoldKNN}{\XX}
\newcommand{\AccFronteira}{\XX}
\newcommand{\TabelaExperimentos}{%
Neurônios ocultos & 2 / 4 / 8 / 16 / 32 / 10-6 & \XX{} $\pm$ \XX & \XX & \XN \\ \midrule
Taxa $\eta$ & 0,005 / 0,01 / 0,05 / 0,1 / 0,5 & \XX{} $\pm$ \XX & \XX & \XN \\ \midrule
Ativação & sigmoid / tanh / ReLU & \XX{} $\pm$ \XX & \XX & \XN \\}
\newcommand{\TabelaEstatisticas}{%
sepal length & \XN & \XN & \XN & \XN & \XN & 0,7826 \\
sepal width  & \XN & \XN & \XN & \XN & \XN & $-$0,4194 \\
petal length & \XN & \XN & \XN & \XN & \XN & 0,9490 \\
petal width  & \XN & \XN & \XN & \XN & \XN & 0,9565 \\}
% Substitui os marcadores pelos valores reais, se o arquivo existir:
\IfFileExists{resultados_macros.tex}{\input{resultados_macros.tex}}{}

\hypersetup{pdftitle={Trabalho 08 - Reconhecimento de Flores Iris usando MLP},
            pdfauthor={Lucas Albino Martins}}

%=====================================================================
\begin{document}

\fazcapa
\fazfolhaderosto

%----------------------------- Resumo --------------------------------
\begin{center}
  {\bfseries RESUMO}
\end{center}
\begin{singlespace}
\noindent
Este trabalho apresenta o projeto, a implementação e a avaliação de uma rede Perceptron
Multicamadas (MLP) para a classificação das três espécies do conjunto de dados Iris (Fisher, 1936)
a partir de quatro medidas morfológicas. A rede foi implementada do zero em NumPy, com função
Softmax numericamente estável na camada de saída, função de custo entropia cruzada categorial e
retropropagação manual, cujos gradientes foram validados por diferenças finitas (erro relativo
máximo de \GCTanh{} para \textit{tanh}). Adotou-se padronização z-score ajustada apenas no treino,
divisão estratificada 70/15/15, gradiente descendente em mini-batch com momentum e parada antecipada.
A arquitetura 4--8--3 (\textit{tanh}) obteve acurácia de \AccTeste{} no conjunto de teste e de
\AccKfold{} $\pm$ \DPKfold{} na validação cruzada estratificada com cinco partições, desempenho
compatível com a literatura e comparável aos baselines de regressão logística multinomial e
\textit{k}-NN. A \textit{Setosa} é reconhecida sem erros, e as confusões restantes concentram-se no
par \textit{Versicolor}--\textit{Virginica}, que não é linearmente separável.

\vspace{0.5cm}
\noindent\textbf{Palavras-chave:} Redes neurais artificiais. Perceptron multicamadas. Softmax.
Entropia cruzada. Retropropagação. Conjunto Iris.
\end{singlespace}
\newpage

\tableofcontents
\newpage

\begin{center}
  {\large\bfseries \TituloTexto\par}
\end{center}
\vspace{0.5cm}

%=====================================================================
\section{Introdução}
\label{sec:introducao}

O conjunto de dados Iris, publicado por Fisher (1936) para
ilustrar a análise discriminante linear, reúne 150 flores de três espécies do gênero \textit{Iris}
--- \textit{I.~setosa}, \textit{I.~versicolor} e \textit{I.~virginica}, com 50 exemplares cada ---
descritas por quatro medidas em centímetros: comprimento e largura da sépala e da pétala. Ele é um
dos \textit{benchmarks} mais citados em reconhecimento de padrões (DUDA; HART, 1973; DASARATHY, 1980; GATES, 1972) porque
reúne, em escala reduzida, uma classe linearmente separável e um par de classes que não é.

Este trabalho atende ao enunciado do Trabalho~08 da disciplina EL056: treinar uma MLP para
classificar as flores, com função de ativação Softmax na saída e função de custo entropia cruzada.
Para além do mínimo exigido, buscou-se rigor metodológico: implementação sem frameworks de
\textit{deep learning}, verificação numérica dos gradientes, ausência de vazamento de dados no
pré-processamento, avaliação por validação cruzada estratificada, estudo de sensibilidade a
hiperparâmetros e comparação com baselines.

O relatório está organizado da seguinte forma. A Seção~\ref{sec:teoria} apresenta a fundamentação
teórica; a Seção~\ref{sec:dados} descreve a base de dados; a Seção~\ref{sec:metodologia}, a
metodologia; as Seções~\ref{sec:resultados} e \ref{sec:discussao} apresentam e discutem os
resultados; e a Seção~\ref{sec:conclusao} conclui o trabalho.

%=====================================================================
\section{Fundamentação Teórica}
\label{sec:teoria}

\subsection{Perceptron Multicamadas e propagação direta}

Uma MLP com $L$ camadas de pesos mapeia uma entrada $\mathbf{x}\in\mathbb{R}^{d}$ em um vetor de
probabilidades $\hat{\mathbf{y}}\in\mathbb{R}^{K}$. Em notação matricial por lote (cada linha de
$\mathbf{A}^{(0)}=\mathbf{X}\in\mathbb{R}^{N\times d}$ é uma amostra):
\begin{align}
  \mathbf{Z}^{(l)} &= \mathbf{A}^{(l-1)}\mathbf{W}^{(l)} + \mathbf{1}\,\mathbf{b}^{(l)\top},
  && l = 1,\dots,L, \label{eq:z}\\
  \mathbf{A}^{(l)} &= f\big(\mathbf{Z}^{(l)}\big), && l = 1,\dots,L-1, \label{eq:a}\\
  \hat{\mathbf{Y}} &= \operatorname{softmax}\big(\mathbf{Z}^{(L)}\big), \label{eq:saida}
\end{align}
em que $\mathbf{W}^{(l)}\in\mathbb{R}^{n_{l-1}\times n_l}$ e $f$ é a ativação das camadas ocultas:
\begin{equation}
  \sigma(z)=\frac{1}{1+e^{-z}},\qquad \tanh(z)=\frac{e^{z}-e^{-z}}{e^{z}+e^{-z}},\qquad
  \operatorname{ReLU}(z)=\max(0,z), \label{eq:ativacoes}
\end{equation}
com derivadas $\sigma'(z)=\sigma(z)[1-\sigma(z)]$, $\tanh'(z)=1-\tanh^2(z)$ e
$\operatorname{ReLU}'(z)=\mathbb{1}[z>0]$ (HAYKIN, 2001).

\subsection{Função Softmax estável}

A Softmax converte os \textit{logits} $\mathbf{z}\in\mathbb{R}^{K}$ em uma distribuição de
probabilidade:
\begin{equation}
  \hat{y}_k=\frac{e^{z_k}}{\sum_{j=1}^{K} e^{z_j}}
  =\frac{e^{z_k-m}}{\sum_{j=1}^{K} e^{z_j-m}},\qquad m=\max_j z_j. \label{eq:softmax}
\end{equation}
A segunda forma é matematicamente idêntica à primeira, pois o fator $e^{-m}$ se cancela no numerador e
no denominador. Numericamente, porém, ela garante $z_k-m\le 0$, de modo que $e^{z_k-m}\in(0,1]$: não há
\textit{overflow} e pelo menos um termo do denominador vale 1.

\subsection{Entropia cruzada categorial}

Com rótulos codificados em \textit{one-hot} ($y_{nk}=1$ se a amostra $n$ pertence à classe $k$), a
função de custo média é
\begin{equation}
  L=-\frac{1}{N}\sum_{n=1}^{N}\sum_{k=1}^{K} y_{nk}\ln \hat{y}_{nk}, \label{eq:ce}
\end{equation}
que corresponde à log-verossimilhança negativa de um modelo multinomial (GOODFELLOW; BENGIO; COURVILLE, 2016). Na
implementação, $\hat{y}$ é limitado ao intervalo $[\varepsilon,1]$, com $\varepsilon=10^{-12}$, para
evitar $\ln 0$.

\subsection{Gradiente conjunto Softmax + entropia cruzada}

O jacobiano da Softmax é $\partial\hat{y}_k/\partial z_j=\hat{y}_k(\delta_{kj}-\hat{y}_j)$. Para uma
amostra:
\begin{equation}
  \frac{\partial L_n}{\partial z_j}
  =-\sum_{k}\frac{y_k}{\hat{y}_k}\,\hat{y}_k(\delta_{kj}-\hat{y}_j)
  =-y_j+\hat{y}_j\sum_{k}y_k
  =\hat{y}_j-y_j, \label{eq:grad_saida}
\end{equation}
pois $\sum_k y_k=1$. Para o custo médio no lote, $\boldsymbol{\delta}^{(L)}=(\hat{\mathbf{Y}}-\mathbf{Y})/N$.
Esse resultado elimina o cálculo explícito do jacobiano, é numericamente estável e tem gradiente
limitado em $[-1,1]$, o que evita o problema de saturação do erro quadrático com saídas sigmoidais.

\subsection{Retropropagação e regra de atualização}

Aplicando a regra da cadeia de trás para frente (RUMELHART; HINTON; WILLIAMS, 1986; HAYKIN, 2001):
\begin{align}
  \nabla_{\mathbf{W}^{(l)}}L &= \mathbf{A}^{(l-1)\top}\boldsymbol{\delta}^{(l)}, \qquad
  \nabla_{\mathbf{b}^{(l)}}L = \boldsymbol{\delta}^{(l)\top}\mathbf{1}, \label{eq:gradW}\\
  \boldsymbol{\delta}^{(l-1)} &= \big(\boldsymbol{\delta}^{(l)}\mathbf{W}^{(l)\top}\big)\odot
  f'\big(\mathbf{Z}^{(l-1)}\big). \label{eq:delta}
\end{align}
Os parâmetros $\theta$ são atualizados por gradiente descendente em mini-batch com momentum:
\begin{equation}
  \mathbf{v}\leftarrow\mu\,\mathbf{v}-\eta\,\nabla_{\theta}L,\qquad
  \theta\leftarrow\theta+\mathbf{v}, \label{eq:momentum}
\end{equation}
em que $\eta$ é a taxa de aprendizado e $\mu$ o coeficiente de momentum; com $\mu=0$ recupera-se o
gradiente descendente puro.

\subsection{Inicialização dos pesos}

Para manter a variância das ativações e dos gradientes aproximadamente constante entre camadas,
adotou-se a inicialização de Xavier/Glorot (GLOROT; BENGIO, 2010), $W\sim\mathcal{N}\big(0,\,2/(n_{in}+n_{out})\big)$,
para \textit{tanh}, sigmoide e a camada Softmax, e a de He (HE \textit{et al.}, 2015), $W\sim\mathcal{N}(0,\,2/n_{in})$,
para camadas ReLU. Os vieses começam em zero.

%=====================================================================
\section{Base de Dados}
\label{sec:dados}

A base foi fornecida como \texttt{iris\_data.xlsx} (150 linhas $\times$ 5 colunas, sem cabeçalho),
com documentação no arquivo \texttt{iris\_names.txt} do repositório UCI (FISHER, 1988). As colunas 1--4
são os atributos em cm (\textit{sepal length}, \textit{sepal width}, \textit{petal length},
\textit{petal width}), armazenados como texto e convertidos para ponto flutuante; a coluna 5 contém
o rótulo. Não há valores ausentes, e as classes são balanceadas (33,3\% cada). As linhas estão
ordenadas por classe em blocos de 50, o que torna obrigatório embaralhar os dados antes de dividi-los.

\subsection{Correção das amostras 35 e 38}

A documentação registra que duas amostras de \textit{Setosa} divergem do artigo original. Na
planilha, ambas valem $(4{,}9;\,3{,}1;\,1{,}5;\,0{,}1)$; os valores de Fisher são
$(4{,}9;\,3{,}1;\,1{,}5;\,\mathbf{0{,}2})$ para a amostra 35 e $(4{,}9;\,\mathbf{3{,}6};\,\mathbf{1{,}4};\,0{,}1)$
para a 38. Optou-se por \textbf{corrigi-las}, por três razões: os valores verdadeiros são conhecidos e
documentados pela própria fonte, de modo que não se trata de imputação; as duas linhas são idênticas
na planilha, uma duplicata espúria que daria peso dobrado a um único ponto; e a correção restaura a
fidelidade aos dados originais. A decisão é parametrizável (\texttt{--corrigir-uci} /
\texttt{CORRIGIR\_UCI}) e tem efeito desprezível sobre o desempenho, pois ambas as amostras pertencem
à classe separável.

\subsection{Análise exploratória}

A Tabela~\ref{tab:estatisticas} compara as estatísticas calculadas (após a correção) com as
documentadas. Sem a correção, os valores reproduzem exatamente a documentação; com ela, apenas
\textit{sepal width} se altera de forma perceptível, por ser o atributo modificado na amostra 38.

\begin{table}[H]
  \centering
  \caption{Estatísticas dos atributos (cm), desvio-padrão amostral e correlação de Pearson com o
  rótulo (0, 1, 2).}
  \label{tab:estatisticas}
  \small
  \begin{tabular}{lcccccc}
    \toprule
    Atributo & Mín. & Máx. & Média & DP & Corr. (calc.) & Corr. (UCI) \\
    \midrule
    \TabelaEstatisticas
    \bottomrule
  \end{tabular}
\end{table}

As Figuras~\ref{fig:scatter} e \ref{fig:boxplots} confirmam o que a documentação aponta: as medidas
de pétala são as mais discriminantes, com correlação com a classe de cerca de 0,95. No plano das
pétalas, a \textit{Setosa} forma um aglomerado isolado, enquanto \textit{Versicolor} e
\textit{Virginica} se sobrepõem parcialmente. A Figura~\ref{fig:correlacao} mostra ainda a forte
colinearidade entre \textit{petal length}, \textit{petal width} e \textit{sepal length}.

\begin{figure}[H]
  \centering
  \includegraphics[width=0.82\textwidth]{scatter_matrix.pdf}
  \caption{Matriz de dispersão dos quatro atributos, colorida por classe (diagonal: histogramas).}
  \label{fig:scatter}
\end{figure}

\begin{figure}[H]
  \centering
  \includegraphics[width=\textwidth]{boxplots.pdf}
  \caption{Distribuição de cada atributo por classe.}
  \label{fig:boxplots}
\end{figure}

\begin{figure}[H]
  \centering
  \includegraphics[width=0.55\textwidth]{correlacao.pdf}
  \caption{Correlação de Pearson entre atributos e rótulo.}
  \label{fig:correlacao}
\end{figure}

%=====================================================================
\section{Metodologia}
\label{sec:metodologia}

\subsection{Pré-processamento e divisão dos dados}

Os rótulos foram codificados em \textit{one-hot}. Para o holdout, adotou-se uma divisão estratificada
70/15/15 com semente fixa (\texttt{SEED = 42}). O arredondamento por classe resulta em 35/8/7
amostras por espécie, ou 105/24/21 no total. O conjunto de validação é usado exclusivamente para a
parada antecipada, e o de teste, apenas na avaliação final. Os atributos foram padronizados por
z-score, $x'=(x-\mu_{\text{treino}})/\sigma_{\text{treino}}$, com média e desvio estimados
\textbf{somente no treino} e aplicados à validação e ao teste, evitando vazamento de informação.

Como 21 amostras de teste produzem uma estimativa de alta variância (cada erro vale cerca de 4,8 pontos
percentuais), o desempenho também foi avaliado por validação cruzada \textbf{k-fold estratificada}
com $k=5$. Em cada partição, 15\% da parte de treino é separada, de forma estratificada, para a parada
antecipada, e o z-score é reajustado sobre os 85\% restantes.

\subsection{Arquitetura e hiperparâmetros}

A Tabela~\ref{tab:config} resume a configuração principal. A escolha de \textbf{uma camada oculta}
apoia-se no teorema da aproximação universal e na simplicidade do problema, cuja única fronteira
difícil (\textit{Versicolor}/\textit{Virginica}) é quase linear. \textbf{Oito neurônios} dão folga
sobre o mínimo necessário sem super-parametrizar as 105 amostras de treino; essa hipótese é testada
na Seção~\ref{sec:experimentos}. A \textbf{tangente hiperbólica} foi escolhida por ser centrada em
zero, o que produz gradientes mais bem condicionados que os da sigmoide, e por ser compatível com a
inicialização de Xavier. O mini-batch de 16 amostras resulta em cerca de 7 atualizações por época, e
o ruído estocástico ajuda a escapar de platôs. O momentum de 0,9 acelera a descida em direções de
curvatura baixa.

\begin{table}[H]
  \centering
  \caption{Configuração principal da MLP.}
  \label{tab:config}
  \small
  \begin{tabular}{ll}
    \toprule
    Item & Valor \\
    \midrule
    Arquitetura & 4 -- 8 (\textit{tanh}) -- 3 (Softmax), \NParam{} parâmetros \\
    Inicialização & Xavier/Glorot (vieses nulos) \\
    Custo & Entropia cruzada categorial, Eq.~\eqref{eq:ce} \\
    Otimizador & SGD em mini-batch, $\eta=0{,}05$, $\mu=0{,}9$, lote = 16 \\
    Parada antecipada & até 2000 épocas, paciência = 100 (perda de validação) \\
    Semente & 42 \\
    \bottomrule
  \end{tabular}
\end{table}

\subsection{Algoritmo de treinamento}

O Algoritmo~\ref{alg:treino} sintetiza o procedimento implementado na função \texttt{treinar}.

\begin{algorithm}[H]
  \caption{Treinamento da MLP com mini-batch, momentum e parada antecipada}
  \label{alg:treino}
  \begin{algorithmic}[1]
    \Require $(\mathbf{X}_{tr},\mathbf{Y}_{tr})$, $(\mathbf{X}_{va},\mathbf{Y}_{va})$, $\eta$, $\mu$,
             lote $B$, $E_{\max}$, paciência $P$
    \Ensure pesos $\theta^\star$ da melhor época
    \State Inicializar $\mathbf{W}^{(l)}$ (Xavier/He), $\mathbf{b}^{(l)}\gets\mathbf{0}$, $\mathbf{v}\gets\mathbf{0}$
    \State $L^\star\gets\infty$; \ $p\gets 0$
    \For{$e = 1$ \textbf{até} $E_{\max}$}
      \State Embaralhar índices de treino
      \For{cada lote $\mathcal{B}$ de tamanho $B$}
        \State $\hat{\mathbf{Y}}_{\mathcal{B}}\gets$ \Call{Forward}{$\mathbf{X}_{\mathcal{B}}$} \Comment{Eqs.~\eqref{eq:z}--\eqref{eq:saida}}
        \State $\boldsymbol{\delta}^{(L)}\gets(\hat{\mathbf{Y}}_{\mathcal{B}}-\mathbf{Y}_{\mathcal{B}})/|\mathcal{B}|$ \Comment{Eq.~\eqref{eq:grad_saida}}
        \State Retropropagar $\boldsymbol{\delta}$ e obter $\nabla_\theta L$ \Comment{Eqs.~\eqref{eq:gradW}--\eqref{eq:delta}}
        \State $\mathbf{v}\gets\mu\mathbf{v}-\eta\nabla_\theta L$; \ $\theta\gets\theta+\mathbf{v}$ \Comment{Eq.~\eqref{eq:momentum}}
      \EndFor
      \State Calcular $L_{va}$ e acurácias de treino/validação; registrar histórico
      \If{$L_{va} < L^\star - 10^{-6}$}
        \State $L^\star\gets L_{va}$; \ $\theta^\star\gets\theta$; \ $p\gets 0$
      \Else
        \State $p\gets p+1$; \ \textbf{se} $p\ge P$ \textbf{então interromper}
      \EndIf
    \EndFor
    \State \Return $\theta^\star$
  \end{algorithmic}
\end{algorithm}

\subsection{Verificação do gradiente}

Antes do treinamento, o gradiente analítico foi comparado ao numérico, obtido por diferenças centrais
($\epsilon=10^{-5}$) em todos os parâmetros de uma rede 4--5--4--3. A rede tem duas camadas ocultas
para exercitar a recursão da Eq.~\eqref{eq:delta}. O critério de aceitação é
$|g_a-g_n|/\max(|g_a|+|g_n|,10^{-12})<10^{-6}$.

\subsection{Experimentos e baselines}
\label{sec:experimentos}

Em torno da configuração base, variou-se um fator por vez, com avaliação por k-fold ($k=5$): número de
neurônios ocultos $\{2, 4, 8, 16, 32, (10,6)\}$, taxa de aprendizado
$\{0{,}005;\ 0{,}01;\ 0{,}05;\ 0{,}1;\ 0{,}5\}$ e ativação $\{\text{sigmoide}, \tanh, \text{ReLU}\}$.
Como baselines, usaram-se a regressão logística multinomial, equivalente a uma rede Softmax sem
camada oculta e, portanto, uma medida direta do ganho trazido pela não linearidade, e o
\textit{k}-NN com $k=5$ (DASARATHY, 1980; GATES, 1972), ambos do scikit-learn, sobre os mesmos splits e
partições. Treinou-se ainda uma MLP auxiliar apenas com os atributos de pétala, para visualizar a
fronteira de decisão em duas dimensões.

\subsection{Métricas}

Para cada classe $c$: precisão $P_c=TP_c/(TP_c+FP_c)$, revocação $R_c=TP_c/(TP_c+FN_c)$ e
$F1_c=2P_cR_c/(P_c+R_c)$, além das médias \textit{macro}, da acurácia global e das matrizes de
confusão.

%=====================================================================
\section{Resultados}
\label{sec:resultados}

Todos os valores desta seção foram gerados por \texttt{iris\_mlp.py} (registrados em
\texttt{resultados.json}). O notebook \texttt{iris\_mlp.ipynb} reproduz os mesmos números.

\subsection{Verificação do gradiente}

O erro relativo máximo entre gradiente analítico e numérico foi \GCSig{} (sigmoide), \GCTanh{}
(\textit{tanh}) e \GCRelu{} (ReLU), todos abaixo do limiar de $10^{-6}$. A implementação da
retropropagação, incluindo a Eq.~\eqref{eq:grad_saida}, está correta.

\subsection{Treinamento}

A Figura~\ref{fig:curvas} mostra as curvas de aprendizado. A parada antecipada selecionou a época
\MelhorEpoca{}, de um total de \EpocasExec{} épocas executadas. Nessa época, a entropia cruzada foi de
\CETreino{} no treino e \CEVal{} na validação, com acurácias de \AccTreino{} e \AccVal{},
respectivamente.

\begin{figure}[H]
  \centering
  \includegraphics[width=\textwidth]{curvas_treino.pdf}
  \caption{Curvas de perda (escala logarítmica) e de acurácia de treino e validação; a linha
  tracejada indica a época restaurada pela parada antecipada.}
  \label{fig:curvas}
\end{figure}

\subsection{Desempenho no conjunto de teste}

No conjunto de teste (21 amostras), a MLP obteve acurácia de \AccTeste{}, F1 macro de \FUmTeste{}
e entropia cruzada de \CETeste{}. As Tabelas~\ref{tab:metricas} e \ref{tab:cm} e a
Figura~\ref{fig:cm} detalham os resultados.

\begin{table}[H]
  \centering
  \caption{Métricas por classe no conjunto de teste.}
  \label{tab:metricas}
  \small
  \begin{tabular}{lccc}
    \toprule
    Classe & Precisão & Revocação & F1 \\
    \midrule
    \textit{Iris setosa}     & \PSet & \RSet & \FSet \\
    \textit{Iris versicolor} & \PVer & \RVer & \FVer \\
    \textit{Iris virginica}  & \PVir & \RVir & \FVir \\
    \midrule
    Média macro & \PMacro & \RMacro & \FMacro \\
    \bottomrule
  \end{tabular}
\end{table}

\begin{table}[H]
  \centering
  \caption{Matrizes de confusão (linhas: classe real; colunas: classe prevista).}
  \label{tab:cm}
  \small
  \begin{tabular}{lccc@{\hspace{1.2cm}}ccc}
    \toprule
    & \multicolumn{3}{c}{Teste (holdout, $n=21$)} & \multicolumn{3}{c}{k-fold agregado ($n=150$)} \\
    \cmidrule(lr){2-4}\cmidrule(lr){5-7}
    Real $\backslash$ Prevista & Set. & Ver. & Vir. & Set. & Ver. & Vir. \\
    \midrule
    Setosa     & \CMTAA & \CMTAB & \CMTAC & \CMKAA & \CMKAB & \CMKAC \\
    Versicolor & \CMTBA & \CMTBB & \CMTBC & \CMKBA & \CMKBB & \CMKBC \\
    Virginica  & \CMTCA & \CMTCB & \CMTCC & \CMKCA & \CMKCB & \CMKCC \\
    \bottomrule
  \end{tabular}
\end{table}

\begin{figure}[H]
  \centering
  \begin{subfigure}{0.48\textwidth}
    \centering
    \includegraphics[width=\linewidth]{matriz_confusao.pdf}
    \caption{Teste (holdout).}
  \end{subfigure}\hfill
  \begin{subfigure}{0.48\textwidth}
    \centering
    \includegraphics[width=\linewidth]{matriz_confusao_kfold.pdf}
    \caption{k-fold agregado (predições fora da partição).}
  \end{subfigure}
  \caption{Matrizes de confusão da MLP 4--8--3.}
  \label{fig:cm}
\end{figure}

\subsection{Validação cruzada e experimentos}

Com $k=5$, a configuração principal obteve acurácia de \AccKfold{} $\pm$ \DPKfold{} e F1 macro de
\FUmKfold{} $\pm$ \DPFUmKfold{}. As acurácias por partição foram \FoldsKfold. A
Tabela~\ref{tab:experimentos} e a Figura~\ref{fig:experimentos} apresentam o estudo de sensibilidade.

\begin{table}[H]
  \centering
  \caption{Sensibilidade a hiperparâmetros (k-fold estratificado, $k=5$; média $\pm$ desvio-padrão).
  Os demais hiperparâmetros seguem a Tabela~\ref{tab:config}.}
  \label{tab:experimentos}
  \small
  \begin{tabular}{llccc}
    \toprule
    Fator & Valor & Acurácia & F1 macro & Melhor época (média) \\
    \midrule
    \TabelaExperimentos
    \bottomrule
  \end{tabular}
\end{table}

\begin{figure}[H]
  \centering
  \includegraphics[width=\textwidth]{experimentos.pdf}
  \caption{Acurácia média do k-fold ($\pm$ desvio-padrão) para cada fator variado.}
  \label{fig:experimentos}
\end{figure}

\subsection{Comparação com baselines}

\begin{table}[H]
  \centering
  \caption{Comparação com baselines (mesmo split holdout e mesmas partições do k-fold).}
  \label{tab:baselines}
  \small
  \begin{tabular}{lcc}
    \toprule
    Modelo & Acurácia (teste) & Acurácia k-fold \\
    \midrule
    MLP 4--8--3 (\textit{tanh}, Softmax + CE) & \AccTeste & \AccKfold{} $\pm$ \DPKfold \\
    Regressão logística multinomial & \AccTesteLR & \AccKfoldLR{} $\pm$ \DPKfoldLR \\
    \textit{k}-NN ($k=5$) & \AccTesteKNN & \AccKfoldKNN{} $\pm$ \DPKfoldKNN \\
    \bottomrule
  \end{tabular}
\end{table}

\subsection{Fronteira de decisão}

A MLP auxiliar treinada apenas com \textit{petal length} e \textit{petal width} obteve acurácia de
teste de \AccFronteira. Suas regiões de decisão aparecem na Figura~\ref{fig:fronteira}.

\begin{figure}[H]
  \centering
  \includegraphics[width=0.75\textwidth]{fronteira_decisao.pdf}
  \caption{Regiões de decisão da MLP no plano das pétalas; as estrelas indicam as amostras de teste.}
  \label{fig:fronteira}
\end{figure}

%=====================================================================
\section{Discussão}
\label{sec:discussao}

\textbf{Correção da implementação.} O \textit{gradient checking} confirma que a retropropagação
manual, incluindo a simplificação $\partial L/\partial\mathbf{z}=\hat{\mathbf{y}}-\mathbf{y}$, está
correta para as três ativações. Esse teste é o que dá confiabilidade a uma implementação feita do
zero: sem ele, erros sutis de sinal ou de dimensão podem passar despercebidos, porque a rede às vezes
"aprende" mesmo com gradientes incorretos.

\textbf{Dinâmica de treinamento.} As curvas da Figura~\ref{fig:curvas} mostram queda rápida da perda
nas primeiras épocas, seguida de estabilização. A perda de validação passa a oscilar em torno do
mínimo, efeito do ruído do mini-batch e do tamanho reduzido do conjunto de validação. A parada
antecipada restaura os pesos de menor perda de validação, o que funciona como regularização
implícita.

\textbf{Holdout versus k-fold.} Com 7 amostras por classe no teste, a acurácia holdout
(\AccTeste) tem alta variância e não deve ser lida isoladamente. A estimativa de referência é a do
k-fold (\AccKfold{} $\pm$ \DPKfold), cujo desvio-padrão quantifica a variabilidade entre partições.

\textbf{Confusão \textit{Versicolor} $\times$ \textit{Virginica}.} A matriz agregada do k-fold
(Tabela~\ref{tab:cm}) mostra que a \textit{Setosa} não é confundida com nenhuma outra classe, o que é
coerente com sua separabilidade linear. Dos \ErrosTotalKfold{} erros fora da partição,
\ErrosVerVir{} ocorrem entre \textit{Versicolor} e \textit{Virginica}. A Figura~\ref{fig:fronteira}
mostra a causa: as duas espécies se sobrepõem na região de pétalas intermediárias (comprimento de
cerca de 4,5 a 5,1~cm e largura de 1,5 a 1,8~cm). Nessa região, amostras de classes distintas têm
medidas praticamente iguais, e o erro é essencialmente irredutível com esses quatro atributos.

\textbf{Sensibilidade a hiperparâmetros.} Na faixa avaliada, as variações de arquitetura, taxa de
aprendizado e ativação produzem diferenças de desempenho da ordem do desvio-padrão entre partições
(Tabela~\ref{tab:experimentos}), ou seja, sem diferença estatisticamente relevante. Taxas menores
exigem mais épocas até o mínimo de validação, e taxas muito altas tornam a otimização mais ruidosa.
Redes com 2 neurônios já são suficientes, o que confirma que a fronteira necessária é simples. A
escolha de 8 neurônios com \textit{tanh} é, portanto, adequada, mas não crítica.

\textbf{MLP versus baselines.} A regressão logística multinomial, que é uma rede Softmax sem camada
oculta, atinge desempenho próximo ao da MLP (Tabela~\ref{tab:baselines}), indicando que o problema é
quase linearmente separável. O ganho possível da camada oculta limita-se às poucas amostras na região
de sobreposição \textit{Versicolor}/\textit{Virginica}. O \textit{k}-NN também apresenta desempenho
elevado, coerente com os resultados históricos para o Iris (DASARATHY, 1980; GATES, 1972).

\textbf{Limitações.} A base é pequena (150 amostras), o que limita o poder estatístico das
comparações entre configurações. Os experimentos variam um fator por vez e não exploram interações.
Por fim, as diferenças observadas entre modelos são menores que a variabilidade entre partições.

%=====================================================================
\section{Conclusão}
\label{sec:conclusao}

Foi implementada do zero, em NumPy, uma MLP para o reconhecimento das três espécies de flores Iris,
atendendo a todos os requisitos do enunciado: três classes balanceadas, quatro atributos em cm,
ativação Softmax (em forma numericamente estável) na saída e função de custo entropia cruzada. A
derivação do gradiente conjunto Softmax + entropia cruzada foi apresentada e validada numericamente.
O protocolo experimental empregado incluiu z-score sem vazamento de dados, divisão estratificada,
parada antecipada, k-fold estratificado, estudo de sensibilidade e baselines.

A rede 4--8--3 obteve \AccTeste{} de acurácia no teste e \AccKfold{} $\pm$ \DPKfold{} no k-fold, com
reconhecimento perfeito da \textit{Setosa} e erros restritos ao par \textit{Versicolor}--\textit{Virginica},
conforme previsto pela análise exploratória. Como trabalhos futuros, sugerem-se testes estatísticos
pareados entre modelos (por exemplo, McNemar ou k-fold repetido), regularização L2 e estimativa de
incerteza das probabilidades Softmax (calibração).

%=====================================================================
\section*{Referências}
\addcontentsline{toc}{section}{Referências}
\begin{singlespace}
\setlength{\parindent}{0pt}
\setlength{\parskip}{0.5\baselineskip}
\raggedright

DASARATHY, B. V. Nosing around the neighborhood: a new system structure and classification rule for
recognition in partially exposed environments. \textbf{IEEE Transactions on Pattern Analysis and
Machine Intelligence}, v. PAMI-2, n. 1, p. 67--71, 1980.

DUDA, R. O.; HART, P. E. \textbf{Pattern classification and scene analysis}. New York: John Wiley \&
Sons, 1973.

FISHER, R. A. The use of multiple measurements in taxonomic problems. \textbf{Annals of Eugenics},
v. 7, n. 2, p. 179--188, 1936.

FISHER, R. A. \textbf{Iris}. UCI Machine Learning Repository, 1988. Disponível em:
\url{http://archive.ics.uci.edu/ml/machine-learning-databases/iris/}. Acesso em: 7 out. 2026.

GATES, G. W. The reduced nearest neighbor rule. \textbf{IEEE Transactions on Information Theory},
v. 18, n. 3, p. 431--433, 1972.

GLOROT, X.; BENGIO, Y. Understanding the difficulty of training deep feedforward neural networks. In:
INTERNATIONAL CONFERENCE ON ARTIFICIAL INTELLIGENCE AND STATISTICS, 13., 2010, Sardinia.
\textbf{Proceedings} [...]. PMLR, 2010. p. 249--256.

GOODFELLOW, I.; BENGIO, Y.; COURVILLE, A. \textbf{Deep learning}. Cambridge, MA: MIT Press,
2016.

HAYKIN, S. \textbf{Redes neurais}: princípios e prática. 2. ed. Porto Alegre: Bookman,
2001.

HE, K.; ZHANG, X.; REN, S.; SUN, J. Delving deep into rectifiers: surpassing human-level performance
on ImageNet classification. In: IEEE INTERNATIONAL CONFERENCE ON COMPUTER VISION, 2015, Santiago.
\textbf{Proceedings} [...]. IEEE, 2015. p. 1026--1034.

RUMELHART, D. E.; HINTON, G. E.; WILLIAMS, R. J. Learning representations by back-propagating errors.
\textbf{Nature}, v. 323, p. 533--536, 1986.
\end{singlespace}

%=====================================================================
\appendix
\section{Trechos do código}
\label{ap:codigo}

Os trechos a seguir foram extraídos de \texttt{iris\_mlp.py} (as docstrings foram omitidas).

\begin{lstlisting}[caption={Softmax numericamente estável (Eq.~\ref{eq:softmax}).},label={lst:softmax}]
def softmax(z: np.ndarray) -> np.ndarray:
    z_desl = z - np.max(z, axis=1, keepdims=True)
    e = np.exp(z_desl)
    return e / np.sum(e, axis=1, keepdims=True)
\end{lstlisting}

\begin{lstlisting}[caption={Entropia cruzada categorial (Eq.~\ref{eq:ce}).},label={lst:ce}]
def cross_entropy(y: np.ndarray, y_hat: np.ndarray, eps: float = 1e-12) -> float:
    return float(-np.mean(np.sum(y * np.log(np.clip(y_hat, eps, 1.0)), axis=1)))
\end{lstlisting}

\begin{lstlisting}[caption={Propagação direta, retropropagação e atualização com momentum (métodos da classe \texttt{MLP}).},label={lst:backprop}]
def forward(self, X: np.ndarray) -> np.ndarray:
    self._A = [X]
    self._Z = []
    a = X
    for l in range(len(self.W)):
        z = a @ self.W[l] + self.b[l]                       # Eq. (1)
        self._Z.append(z)
        a = softmax(z) if l == len(self.W) - 1 else self.f(z)
        self._A.append(a)
    return a

def backward(self, Y: np.ndarray):
    N = Y.shape[0]
    delta = (self._A[-1] - Y) / N                           # Softmax + CE
    dW = [np.empty(0)] * len(self.W)
    db = [np.empty(0)] * len(self.W)
    for l in range(len(self.W) - 1, -1, -1):
        dW[l] = self._A[l].T @ delta                        # gradiente dos pesos
        db[l] = delta.sum(axis=0)
        if l > 0:
            delta = (delta @ self.W[l].T) * self.df(self._Z[l - 1])
    return dW, db

def update(self, dW, db, lr: float, momentum: float = 0.0) -> None:
    for l in range(len(self.W)):
        self.vW[l] = momentum * self.vW[l] - lr * dW[l]
        self.vb[l] = momentum * self.vb[l] - lr * db[l]
        self.W[l] += self.vW[l]
        self.b[l] += self.vb[l]
\end{lstlisting}

\end{document}
